\sectionheading{5.14 ইন্টারপ্রেটার (Interpreter)}

High Level Language-এ লেখা Program কম্পিউটার সরাসরি বুঝতে পারে না।
\textbf{Interpreter} Source Program-এর Instruction সাধারণত
\textbf{একটি করে পড়ে, অনুবাদ করে এবং Execute করে}।

একটি Instruction-এর কাজ শেষ হলে পরবর্তী Instruction-এ যায়।
কোনো Error পাওয়া গেলে সাধারণত সেই স্থানেই কাজ থেমে যায়।
এই কারণে Program পরীক্ষা ও ভুল সংশোধনের সময় Interpreter
ব্যবহার সুবিধাজনক।

%================================================
% Definition
%================================================
\booksubheading{Interpreter কাকে বলে?}

\begin{definitionbox}

\textbf{\color{darkTeal}সংজ্ঞা}

যে Translator Program High Level Language-এ লেখা Source Program-কে
\textbf{একটি করে Statement বা Line অনুবাদ করে এবং সঙ্গে সঙ্গে
Execute করে}, তাকে \textbf{Interpreter} বলে।

\end{definitionbox}

\textbf{সহজভাবে মনে রাখবে—}

\begin{center}

\textbf{
একটি Line পড়া
\(\longrightarrow\)
অনুবাদ
\(\longrightarrow\)
Execute
\(\longrightarrow\)
পরের Line
}

\end{center}

%================================================
% Working
%================================================
\booksubheading{Interpreter কীভাবে কাজ করে?}

Interpreter সম্পূর্ণ Program আগে অনুবাদ না করে সাধারণত
একটি Statement নিয়ে কাজ করে।

\begin{center}
\begin{adjustbox}{max width=0.98\linewidth}
\fcolorbox{darkBlue}{softBlue}{
    \parbox[c][9mm][c]{0.20\linewidth}{
        \centering\textbf{Source\\Statement}
    }
}
\quad
\(\longrightarrow\)
\quad
\fcolorbox{mainTeal}{softGreen}{
    \parbox[c][9mm][c]{0.18\linewidth}{
        \centering\textbf{Interpreter}
    }
}
\quad
\(\longrightarrow\)
\quad
\fcolorbox{darkBlue}{softBlue}{
    \parbox[c][9mm][c]{0.20\linewidth}{
        \centering\textbf{Translation}
    }
}
\quad
\(\longrightarrow\)
\quad
\fcolorbox{mainTeal}{softGreen}{
    \parbox[c][9mm][c]{0.17\linewidth}{
        \centering\textbf{Execution}
    }
}
\end{adjustbox}
\end{center}

এরপর Interpreter পরবর্তী Statement-এ যায়।

%================================================
% Example
%================================================
\begin{examplebox}{সহজ উদাহরণ}

ধরা যাক একটি Program-এ তিনটি Statement আছে—

\begin{center}

\textenglish{\ttfamily Statement 1}

\textenglish{\ttfamily Statement 2}

\textenglish{\ttfamily Statement 3}

\end{center}

Interpreter প্রথমে \textenglish{\ttfamily Statement 1}
অনুবাদ ও Execute করবে।

এরপর \textenglish{\ttfamily Statement 2}, তারপর
\textenglish{\ttfamily Statement 3}-এ যাবে।

যদি \textenglish{\ttfamily Statement 2}-এ Error থাকে,
তাহলে সাধারণত সেখানেই Execution থেমে যাবে।

\end{examplebox}

%================================================
% Features
%================================================
\booksubheading{Interpreter-এর প্রধান বৈশিষ্ট্য}

\begin{tabularx}{\linewidth}{|>{\bfseries}p{0.28\linewidth}|Y|}
\hline

\rowcolor{softBlue}
\multicolumn{1}{|c|}{\textbf{বৈশিষ্ট্য}}
&
\multicolumn{1}{c|}{\textbf{ব্যাখ্যা}}
\\
\hline

Translator Program
&
High Level Program অনুবাদে ব্যবহৃত হয়।
\\
\hline

Statement ভিত্তিক
&
সাধারণত একটি করে Statement বা Line অনুবাদ করে।
\\
\hline

তাৎক্ষণিক Execution
&
অনুবাদের পর Instruction সঙ্গে সঙ্গে Execute করে।
\\
\hline

Error-এ থামে
&
Error পাওয়া গেলে সাধারণত সেই Statement-এ Execution বন্ধ হয়।
\\
\hline

Object Program
&
সাধারণত পৃথক পূর্ণ Object Program সংরক্ষণ করে না।
\\
\hline

Debugging সহজ
&
Error-এর স্থান দ্রুত শনাক্ত করা সুবিধাজনক।
\\
\hline

পুনরায় Translation
&
Program আবার চালালে Statementগুলো পুনরায় অনুবাদ করতে হয়।
\\
\hline

\end{tabularx}

%================================================
% Error Detection
%================================================
\booksubheading{Interpreter-এ Error শনাক্তকরণ}

Interpreter একটি Statement অনুবাদ ও Execute করার পর পরবর্তী
Statement-এ যায়।

Error না থাকলে—

\begin{center}

\textbf{
Line পড়া
\(\longrightarrow\)
অনুবাদ
\(\longrightarrow\)
Execute
\(\longrightarrow\)
Next Line
}

\end{center}

কিন্তু Error থাকলে—

\begin{center}

\textbf{
Line পড়া
\(\longrightarrow\)
Error
\(\longrightarrow\)
Execution বন্ধ
\(\longrightarrow\)
Error সংশোধন
}

\end{center}

এই কারণে Interpreter-এ Error-এর অবস্থান দ্রুত চিহ্নিত করা
তুলনামূলক সহজ।

%================================================
% Advantages
%================================================
\booksubheading{Interpreter-এর সুবিধা}

\begin{compacttwocolpoints}

\item একটি করে Instruction নিয়ে কাজ করায় Error শনাক্ত করা সহজ।

\item Program পরিবর্তন ও সংশোধন তুলনামূলক সহজ।

\item Debugging ও Testing দ্রুত করা যায়।

\item সাধারণত পৃথক পূর্ণ Object Program সংরক্ষণ না করায়
অতিরিক্ত Storage-এর প্রয়োজন কম হতে পারে।

\item Interactive Programming ও শেখার জন্য সুবিধাজনক।

\item ছোট Program দ্রুত পরীক্ষা করা যায়।

\item Interpreter-কে তুলনামূলক
\textbf{User Friendly} হিসেবে বিবেচনা করা হয়।

\end{compacttwocolpoints}

%================================================
% Disadvantages
%================================================
\booksubheading{Interpreter-এর অসুবিধা}

\begin{itemize}

\item Program চালানোর সময় Instruction বারবার অনুবাদ করতে হয়।

\item Compiled Program-এর তুলনায় Execution ধীর হতে পারে।

\item Program যতবার চালানো হয়, সাধারণত ততবার প্রয়োজনীয়
Translation করতে হয়।

\item বড় Program-এর ক্ষেত্রে Execution-এর সময় বেশি লাগতে পারে।

\end{itemize}

%================================================
% Extra Fact
%================================================
\begin{rememberbox}

\textbf{\color{deepOrange}অতিরিক্ত MCQ তথ্য}

প্রদত্ত পাঠ্যসূত্রে উল্লেখ করা হয়েছে, Compiler-এর তুলনায়
Interpreter-এ Program চালাতে প্রায়
\textbf{৫ থেকে ২৫ গুণ বেশি সময় লাগতে পারে}।

এটি মূল ধারণার তুলনায় কম গুরুত্বপূর্ণ; MCQ-এর জন্য সংক্ষেপে
মনে রাখলেই যথেষ্ট।

\end{rememberbox}

%================================================
% Examples
%================================================
\booksubheading{Interpreter-এর উদাহরণ}

Interpreter-ভিত্তিক Programming Environment-এর পরিচিত উদাহরণ
হিসেবে \textbf{Python} উল্লেখ করা হয়।

প্রদত্ত পাঠ্যসূত্রে Interpreter-এর উদাহরণের ক্ষেত্রে
\textbf{Java ও Python}-এর নাম উল্লেখ রয়েছে।

%================================================
% Compiler vs Interpreter
%================================================
\booksubheading{Compiler ও Interpreter-এর পার্থক্য}

\begin{tabularx}{\linewidth}{|>{\bfseries}p{0.22\linewidth}|Y|Y|}
\hline

\rowcolor{softBlue}
\multicolumn{1}{|c|}{\textbf{বিষয়}}
&
\multicolumn{1}{c|}{\textbf{Compiler}}
&
\multicolumn{1}{c|}{\textbf{Interpreter}}
\\
\hline

অনুবাদের ধরন
&
সম্পূর্ণ Program একসঙ্গে অনুবাদ করে
&
সাধারণত একটি Statement করে অনুবাদ করে
\\
\hline

Error প্রদর্শন
&
Program বিশ্লেষণ শেষে একাধিক Error দেখাতে পারে
&
Error পাওয়া Statement-এ সাধারণত থামে
\\
\hline

Debugging
&
তুলনামূলক কম সরাসরি
&
তুলনামূলক সহজ
\\
\hline

Execution
&
Compile-এর পর তুলনামূলক দ্রুত হতে পারে
&
অনুবাদসহ চলায় তুলনামূলক ধীর
\\
\hline

Object Program
&
সাধারণত তৈরি করে
&
সাধারণত পৃথক পূর্ণ Object Program তৈরি করে না
\\
\hline

পুনরায় অনুবাদ
&
Source অপরিবর্তিত থাকলে প্রতিবার পুরো Compilation প্রয়োজন নাও হতে পারে
&
Program চালানোর সময় পুনরায় Translation প্রয়োজন হয়
\\
\hline

Memory
&
Object Program সংরক্ষণের জন্য বেশি Memory লাগতে পারে
&
পৃথক Object Program না থাকায় তুলনামূলক কম Memory লাগতে পারে
\\
\hline

কাজের প্রক্রিয়া
&
তুলনামূলক জটিল
&
তুলনামূলক সহজ ও User Friendly
\\
\hline

\end{tabularx}

%================================================
% Shortcut
%================================================
\begin{rememberbox}

\textbf{\color{deepOrange}দ্রুত মনে রাখার কৌশল}

\begin{center}

\textbf{
Compiler
\(\rightarrow\)
Whole Program
\(\rightarrow\)
Object Code
\(\rightarrow\)
Faster Execution
}

\vspace{0.5mm}

\textbf{
Interpreter
\(\rightarrow\)
Line by Line
\(\rightarrow\)
Immediate Execution
\(\rightarrow\)
Easier Debugging
}

\end{center}

\end{rememberbox}

%================================================
% Final Reminder
%================================================
\begin{rememberbox}

\textbf{\color{deepOrange}মনে রাখবে}

\begin{itemize}

\item Interpreter একটি \textbf{Translator Program}।

\item এটি High Level Program-এর Instruction সাধারণত
\textbf{একটি করে অনুবাদ ও Execute} করে।

\item Error পাওয়া গেলে সাধারণত
\textbf{সেই Statement-এ থেমে যায়}।

\item Interpreter সাধারণত পৃথক পূর্ণ
\textbf{Object Program তৈরি করে না}।

\item Program আবার চালালে প্রয়োজনীয় Statementগুলো পুনরায়
অনুবাদ করতে হয়।

\item Interpreter-এ \textbf{Debugging সহজ}, কিন্তু Execution
তুলনামূলক ধীর হতে পারে।

\item \textbf{Python} Interpreter-ভিত্তিক Programming-এর
পরিচিত উদাহরণ।

\item সবচেয়ে গুরুত্বপূর্ণ পার্থক্য—

\begin{center}

\textbf{
Compiler = পুরো Program একসঙ্গে
}

\textbf{
Interpreter = Statement ধরে ধরে
}

\end{center}

\end{itemize}

\end{rememberbox}

\small

%================================================
% 5.15 Program Organization
%================================================
